% (C) 2026 Ralf Stephan, in collaboration with Claude Code.
% Released under CC0 1.0 Universal (public-domain dedication).
% See https://creativecommons.org/publicdomain/zero/1.0/

\documentclass[11pt]{article}
\usepackage[T1]{fontenc}
\usepackage[utf8]{inputenc}
\usepackage{lmodern}
\usepackage[margin=2.5cm]{geometry}
\usepackage{amsmath,amssymb}
\usepackage{booktabs,xltabular}
\usepackage{braket}
\usepackage[hidelinks]{hyperref}

% Lean names and paths: verbatim, breakable at . _ /
\DeclareUrlCommand\lean{\urlstyle{tt}}
\def\UrlBreaks{\do\.\do\_\do\/}
\newcolumntype{L}{>{\raggedright\arraybackslash}X}
\newcolumntype{B}{>{\raggedright\arraybackslash}p{0.42\textwidth}}
\emergencystretch=6em

\title{Lean formalization of an Assumptions of Physics technical brief:\\ log}
\author{}
\date{}

\begin{document}
\maketitle

\noindent
This log records which statements of the technical brief in this directory were checked in Lean,
where the formal version differs from the brief, and any problems found in the brief.
The brief is available at \url{https://assumptionsofphysics.org/resources/briefs}.

\section*{Setup}

\begin{itemize}
\item Lean \texttt{v4.36.0-rc2}; Mathlib at commit \texttt{1b1ccf4b}; physlib (\texttt{master})
  and TauCeti (commit \texttt{e7e7f52}, on the same Mathlib) as further dependencies.
\item Build: \texttt{lake build} in the project root.
\item Every file is checked with \texttt{\#print axioms} on its main results. A result is
  recorded here only if it depends on no axioms beyond \lean{propext},
  \lean{Classical.choice} and \lean{Quot.sound}, and so has no \lean{sorry}.
\item Imports are kept to the Mathlib modules actually needed.
\item Material that belongs in Mathlib is kept in \lean{ForMathlib/}, in Mathlib's directory
  layout (library \lean{ForMathlib}), and imported by the brief files.
\end{itemize}


%======================================================================

\section{Physical vector spaces for quantum mechanics}

\paragraph{Files.}\leavevmode
\begin{itemize}
\item \raggedright Files: \lean{briefs/QuantumVectorSpaces/QuantumVectorSpaces.lean}, and\newline
  \lean{briefs/QuantumVectorSpaces/QuantumVectorSpaces/Countable.lean} (a countable generating
  set of observables, Sec.~5), and\newline
  \lean{briefs/QuantumVectorSpaces/QuantumVectorSpaces/PureStates.lean} (second countability
  from the pure states, Sec.~5), and\newline
  \lean{briefs/QuantumVectorSpaces/QuantumVectorSpaces/Frames.lean} (frame changes, Theorem~5,
  [B]WDQ), and\newline
  \lean{briefs/QuantumVectorSpaces/QuantumVectorSpaces/Isomorphisms.lean} (isomorphisms of
  quantum vector spaces, Example~3 and Theorem~5), and\newline
  \lean{briefs/QuantumVectorSpaces/QuantumVectorSpaces/SchwartzIsomorphisms.lean}
  ($\mathcal{S}(\mathbb{R}^n)$ and $\mathcal{S}(\mathbb{R}^m)$ are Lie isomorphic iff $n = m$,
  Example~3)
\item Library: \lean{QuantumVectorSpaces}
\item \raggedright Imports: \lean{SchwartzNorms}
  (the $L^2$ seminorm family on $\mathcal{S}(\mathbb{R})$)
  and new \lean{ForMathlib} files:\newline
  \lean{ForMathlib/Analysis/InnerProductSpace/Defs.lean} (the seminorm
  $\sqrt{\operatorname{re}\langle x, x\rangle}$ of a \lean{PreInnerProductSpace.Core}, and
  continuity of the inner product against it),\newline
  \lean{ForMathlib/Analysis/CStarAlgebra/Matrix.lean} (\lean{Matrix.toEuclideanLin} is
  multiplicative),\newline
  \lean{ForMathlib/Analysis/Distribution/SchwartzSpace/Complete.lean}
  (\lean{SchwartzMap.instCompleteSpace}: $\mathcal{S}(E, F)$ is complete for complete $F$),\newline
  \lean{ForMathlib/Analysis/Distribution/SchwartzSpace/SupNorm.lean}
  (\lean{SchwartzMap.exists_norm_le_toLp_two}: a Sobolev bound for the values of a Schwartz
  function, from TauCeti's \lean{SchwartzMap.exists_integral_norm_fourierInv_le}),\newline
  \lean{ForMathlib/Analysis/InnerProductSpace/MultilinearBasis.lean}\newline
  (\lean{ContinuousMultilinearMap.norm_le_sum_orthonormalBasis}: a multilinear map is bounded
  by its values on an orthonormal basis),\newline
  \lean{ForMathlib/Analysis/LocallyConvex/FirstCountable.lean}
  (\lean{WithSeminorms.exists_countable}: a family of seminorms inducing a first-countable
  topology has a countable subfamily inducing it),\newline
  \lean{ForMathlib/Topology/Algebra/Projectivization/SecondCountable.lean}
  (\lean{Projectivization.secondCountableTopology_of}: with a separating dual, $V$ is second
  countable if its projective space is),\newline
  \lean{ForMathlib/Algebra/Lie/Subalgebra.lean}
  (\lean{LieSubalgebra.coe_lieSpan_eq_span_of_forall_lie_mem}: a set whose brackets lie in its
  span generates its span as a Lie subalgebra)
\end{itemize}

\paragraph{Setting.} $V$ is a complex vector space with an inner product given as a
\lean{PreInnerProductSpace.Core}. Unlike Mathlib's \lean{InnerProductSpace}, this puts no norm
or topology on $V$. The defining observables are a family $X_i$ ($i \in \iota$) in
\lean{Module.End}, and \lean{algebra} is the unital algebra they generate
(\lean{Algebra.adjoin}). The seminorms
$p_A(v) = \sqrt{\langle Av, Av\rangle}$ form \lean{seminormFamily}, and \lean{topology}
is the topology they generate. A topology is \lean{IsAdmissible} if it makes $V$ a topological
vector space and makes the $X_i$ and the inner product (jointly) continuous. For Example~3,
$X f = x f$ and $P f = -i\hbar f'$ act on Mathlib's $\mathcal{S}(\mathbb{R}, \mathbb{C})$
(\lean{posOp}, \lean{momOp}), with the $L^2$ inner product \lean{l2Core}. For Example~4,
$S_i = (\hbar/2)\sigma_i$ acts on $\mathbb{C}^2$ as \lean{EuclideanSpace} (\lean{spin}).
In $n$ dimensions, $X^i f = x_i f$ and $P_i f = -i\hbar\, \partial_{e_i} f$ act on
$\mathcal{S}(\mathbb{R}^n, \mathbb{C})$, with $\mathbb{R}^n$ as \lean{EuclideanSpace ℝ (Fin n)}
(\lean{posOpN}, \lean{momOpN}, indexed by \lean{ObsN n}), and the $L^2$ inner product
\lean{l2CoreN}.

Definition~1 is \lean{IsQuantumVectorSpace c X ħ C}: the $X_i$ are symmetric, they satisfy
$[X_i, X_j] = i\hbar \sum_k C_{ijk} X_k$ with finitely many nonzero $C_{ijk}$ for each $i, j$,
and $V$ with \lean{topology} is complete and second-countable. Completeness refers to
\lean{uniformSpace}, the uniform structure that every topological vector space carries.

\paragraph{Results.}
\begin{xltabular}{\textwidth}{@{}BL@{}}
\toprule
Brief & Lean \\
\midrule
\endhead
Corollary 2, coarsest topology & \lean{isGreatest_topology} (from \lean{isAdmissible_topology}
  and \lean{IsAdmissible.continuous_algebra}) \\
Corollary 2, locally convex & \lean{locallyConvexSpace} \\
finite dimension (Example 4) & \lean{topology_eq_moduleTopology}, \lean{topology_qubit} \\
Example 3, $[X, P] = i\hbar$ & \lean{posOp_mul_momOp_sub} \\
Example 3, $X$, $P$, $I$ symmetric & \lean{isSymmetric_obs} \\
Example 3, Remark: the topology is that of $\mathcal{S}$ & \lean{topology_schwartz},
  \lean{withSeminorms_schwartz}, \lean{isAdmissible_schwartz} (for $\hbar \ne 0$) \\
Example 4, $[S_i, S_j] = i\hbar\,\epsilon_{ijk} S_k$ & \lean{spin_comm} \\
Example 4, $S_i$ symmetric & \lean{isSymmetric_spin} \\
Example 4 is a quantum vector space & \lean{isQuantumVectorSpace_qubit} \\
Example 3, second countability & \lean{secondCountableTopology_schwartz} \\
Example 3, completeness & \lean{SchwartzMap.instCompleteSpace} \\
Example 3 is a quantum vector space & \lean{isQuantumVectorSpace_schwartz} (for
  $\hbar \ne 0$) \\
Example 3 in $n$ dimensions, commutation relations & \lean{posOpN_comm},
  \lean{momOpN_comm}, \lean{posOpN_comm_momOpN} \\
Example 3 in $n$ dimensions, symmetric observables & \lean{isSymmetric_obsN} \\
Example 3 in $n$ dimensions, Remark: the topology is that of $\mathcal{S}$ &
  \lean{topology_schwartzN} (for $\hbar \ne 0$) \\
Example 3 in $n$ dimensions, second countability & \lean{secondCountableTopology_schwartzN} \\
Example 3 in $n$ dimensions is a quantum vector space & \lean{isQuantumVectorSpace_schwartzN}
  (for $\hbar \ne 0$) \\
Theorem 5, [B]COBS & \lean{IsAdmissible.continuous_expectation} \\
Sec.~5, a countable subset of observables generates the topology &
  \lean{exists_countable_topology_eq}, \lean{exists_countable_topology_eq_iff} (under first
  countability) \\
Sec.~5, the same for a quantum vector space & \lean{IsQuantumVectorSpace.exists_countable} \\
Sec.~5, separability is not enough & \lean{Diagonal.separableSpace},
  \lean{Diagonal.not_firstCountableTopology}, \lean{Diagonal.not_exists_countable} \\
Sec.~5, second countability from the pure states &
  \lean{secondCountableTopology_iff_projectivization} (for a definite inner product) \\
Sec.~5, the pure states of a quantum vector space are second countable &
  \lean{IsQuantumVectorSpace.secondCountableTopology_projectivization} \\
Theorem 5, [B]WDQ: two frames with each other's defining observables as observables have the
  same observables & \lean{topology_eq_iff_continuous}, \lean{isObservable_iff_of_frames} \\
Theorem 5, [B]WDQ: sufficient conditions & \lean{topology_eq_of_mem_algebra},
  \lean{topology_eq_of_finiteDimensional} \\
Theorem 5, [B]WDQ: the premise is not automatic & \lean{Diagonal.not_continuous_mulOp} \\
Isomorphisms of quantum vector spaces & \lean{UnitaryEquiv}, \lean{LieIso} (with
  \lean{refl}, \lean{symm}, \lean{trans}), \lean{LieIso.toUnitaryEquiv} \\
Isomorphisms preserve observables and pure states & \lean{UnitaryEquiv.isObservable_conj_iff},
  \lean{UnitaryEquiv.pureStates} \\
Isomorphisms and frame changes & \lean{induced_topology_eq_pullback},
  \lean{induced_topology_eq_iff} \\
Example 3, $\mathcal{S}(\mathbb{R}^n)$ and $\mathcal{S}(\mathbb{R}^m)$ isomorphic iff $n = m$ &
  \lean{nonempty_lieIso_schwartzN_iff} (for Lie isomorphisms and $\hbar, \hbar' \ne 0$; from
  \lean{LieIso.finrank_lieAlgebra}, \lean{linearIndependent_obsN},
  \lean{finrank_lieAlgebra_obsN}) \\
\bottomrule
\end{xltabular}

\paragraph{Differences from the brief.}
\begin{itemize}
\item ``Coarsest'' is \lean{IsGreatest} in Mathlib's order on topologies, where coarser is
  larger.
\item The inner product is only assumed positive semidefinite in Corollary~2; definiteness is
  used only in finite dimension, to make the topology Hausdorff.
\item Example~3 is formalized twice. For $\mathcal{S}(\mathbb{R})$, the proof that the
  topologies agree has two halves. (a) Each seminorm $\|x^a \partial^b f\|_2$
  equals $|\hbar|^{-b} p_{X^a P^b}(f)$. (b) Every element of the algebra is a combination of the
  normally ordered monomials $X^a P^b$ (\lean{algebra_le_monoSpan}, from
  $P X^{n+1} = X^{n+1} P - (n+1) i\hbar X^n$), so every $p_A$ is bounded by finitely many
  $\|x^a \partial^b f\|_2$. With \lean{withSeminorms_l1_l2} from the previous brief, this gives
  the Schwartz topology.
\item For $\mathcal{S}(\mathbb{R}^n)$, one half is that the Schwartz topology is admissible
  (\lean{isAdmissible_schwartzN}), hence finer than the topology of Corollary~2. For the other
  half, each Schwartz seminorm $\sup_x \|x\|^k \|D^m f(x)\|$ is bounded by finitely many $p_A$
  (\lean{exists_seminorm_le}). The derivative $D^m f(x)$ is bounded by its values on the unit
  vectors, $\|x\|^k \le (1 + \|x\|^2)^k$ is the polynomial $(1 + \sum_i (X^i)^2)^k$, and the
  Sobolev bound $|g(x)| \le C (\|g\|_2 + \sum_i \|\partial_i^{n+1} g\|_2)$ turns the supremum
  into $L^2$ norms of elements of the algebra applied to $f$. No normal ordering is needed, since
  the operators $\partial_i^{n+1} (1 + \sum_j (X^j)^2)^k \partial^r$ are themselves in the
  algebra.
\item In finite dimension the topology of Corollary~2 is the unique Hausdorff vector-space
  topology, whatever the defining observables (\lean{topology_eq_moduleTopology}). This is how
  Example~4 is checked.
\item Second countability in $n$ dimensions comes from a general fact
  (\lean{secondCountableTopology_topology}). With countably many defining observables, the $p_A$
  for the words $A = X_{i_1} \cdots X_{i_k}$ already generate the topology, so $v \mapsto (Av)_A$
  embeds $V$ into a countable product of copies of $L^2$.
\item Second countability of $\mathcal{S}(\mathbb{R})$ is proved by embedding it into a countable
  product of copies of $L^2(\mathbb{R})$ through $f \mapsto (x^a \partial^b f)_{a,b}$
  (\lean{l2Embedding}). The topology of $\mathcal{S}(\mathbb{R})$ is induced by this map,
  because it is given by the $L^2$ seminorms (\lean{withSeminorms_l1_l2}).
\item Completeness of $\mathcal{S}(E, F)$, for any complete $F$, is proved in
  \lean{ForMathlib}, since neither Mathlib (at this pin or on master) nor TauCeti nor physlib
  has it. The proof is the classical one (Becnel and Sengupta, \emph{Mathematics} 3 (2015),
  \S3.5): along a Cauchy sequence the weighted derivatives $\|x\|^k D^n u_m$ converge
  uniformly; the limits are the derivatives of the limit function, by the theorem on uniform
  limits of derivatives (\lean{hasFDerivAt_of_tendstoUniformly}); and the bounds pass to the
  limit. For the qubit, completeness and second countability come from finite dimension.
\item The claim of Sec.~5 that a countable subset of the observables generates the topology is
  proved under first countability, which second countability implies; it is in fact equivalent
  to first countability (\lean{exists_countable_topology_eq_iff}). The proof
  (\lean{WithSeminorms.exists_countable}) takes a countable basis of balls at $0$; each ball
  involves finitely many words $X_{i_1} \cdots X_{i_k}$, and their letters form the countable
  subset. For a quantum vector space the subset is enlarged to be closed under the structure
  constants, so that it satisfies Definition~1 itself
  (\lean{IsQuantumVectorSpace.exists_countable}).
\item Separability is not enough. On the finitely supported sequences, with the $\ell^2$ inner
  product and every multiplication operator $(M_f v)_n = f_n v_n$ by a real sequence $f$ as a
  defining observable, the observables are symmetric and commute (\lean{Diagonal.isSymmetric_mulOp},
  \lean{Diagonal.mulOp_comm}), and the space is separable, as the union of countably many
  finite-dimensional subspaces (\lean{Diagonal.separableSpace}). But the topology is not first
  countable: for countably many balls $U_n$ at $0$, choose $t_n > 0$ with $t_n e_n \in U_n$; then
  $\|M_f(t_n e_n)\| = 1$ for $f_n = 1/t_n$, so no $U_n$ lies in the ball of $p_{M_f}$
  (\lean{Diagonal.not_firstCountableTopology}). Hence no countable subset of the observables
  generates the topology (\lean{Diagonal.not_exists_countable}). The topology is the finest
  locally convex one, so the space is also complete; this is not formalized.
\item The pure states are the rays of $V$, that is, the projective space $\mathbb{P}(V)$. Its
  topology is the quotient topology from the nonzero vectors, which is Mathlib's
  \lean{Projectivization.instTopologicalSpace}. With a definite inner product, $V$ is second
  countable if and only if $\mathbb{P}(V)$ is
  (\lean{secondCountableTopology_iff_projectivization}), so the claim of Sec.~5 holds with the
  topology of Definition~1 itself; nothing has to be extended. One direction holds because the
  quotient map is open. For the other, on the open set where $\langle v_0, v\rangle \ne 0$ the
  map $v \mapsto ([v], \langle v_0, v\rangle)$ is an embedding into the second-countable
  space $\mathbb{P}(V) \times \mathbb{C}$ (\lean{Projectivization.isInducing_chart}); countably
  many such sets cover the nonzero vectors, because $\mathbb{P}(V)$ is Lindel\"of, and a translate
  of one of them covers $0$. The topology that the ensembles would induce on the pure states is
  not compared with the quotient topology, since the brief defines no ensemble space over $V$.
\item For [B]WDQ, two frames are two families of defining observables $X^A_i$ and $X^B_j$ on the
  same space $V$ with the same inner product, and an observable of a frame is a symmetric
  operator that is continuous for its topology (\lean{IsObservable}). The brief's premise, that
  the defining observables of each frame are observables in the other, is a physical assumption
  and is kept as a hypothesis. It holds if and only if the two frames give the same topology
  (\lean{topology_eq_iff_continuous}): the topology of $B$ is then admissible for $A$, hence
  finer than that of $A$ by Corollary~2, and vice versa. The two frames then have the same
  observables (\lean{isObservable_iff_of_frames}), so the identity is a linear homeomorphism that
  preserves the inner product and all observables. The premise holds when the defining
  observables of each frame lie in the algebra of the other (\lean{topology_eq_of_mem_algebra}),
  and always in finite dimension (\lean{topology_eq_of_finiteDimensional}). It does not follow
  from sharing the space and the inner product: on the finitely supported sequences, the
  multiplication by $n$, a defining observable of the frame with all multiplication operators,
  is not continuous for the frame whose only defining observable is the identity
  (\lean{Diagonal.not_continuous_mulOp}), since $\|e_n\| = 1$ while $\|n e_n\| = n$.
\item The brief does not define isomorphisms of quantum vector spaces. Two notions are
  formalized, each closed under identity, inverse and composition. A \emph{unitary equivalence}
  (\lean{UnitaryEquiv}) is a linear bijection $U\colon V \to W$ that preserves the inner product
  and carries the topology of $W$ to that of $V$. Conjugation $A \mapsto U A U^{-1}$ is then a
  bijection of the observables (\lean{UnitaryEquiv.isObservable_conj_iff}), $U$ induces a
  homeomorphism of the pure states (\lean{UnitaryEquiv.pureStates}), and the Born rule is
  preserved with the inner product. Along any $U$ that preserves the inner product, the topology
  of $W$ pulls back to the topology on $V$ with defining observables $U^{-1} Y_j U$
  (\lean{induced_topology_eq_pullback}). So $U$ is a unitary equivalence if and only if the
  $U^{-1} Y_j U$ are continuous on $V$ and the $U X_i U^{-1}$ are continuous on $W$
  (\lean{induced_topology_eq_iff}); for $U = \mathrm{id}$ this is the frame condition above. A
  \emph{Lie isomorphism} (\lean{LieIso}) is a linear bijection that preserves the inner product
  and whose conjugation maps the complex Lie algebra spanned by the $X_i$ (\lean{lieAlgebra}),
  with the commutator as bracket, onto that spanned by the $Y_j$. This is the brief's
  ``preserve the Lie algebraic structure''. Every Lie isomorphism is a unitary equivalence
  (\lean{LieIso.toUnitaryEquiv}): the two Lie algebras generate conjugate associative algebras,
  hence the same topology.
\item The truth of the ``iff $n = m$'' part of Example~3 depends on the notion of isomorphism.
  For Lie isomorphisms it holds (\lean{nonempty_lieIso_schwartzN_iff}), even with different
  nonzero values $\hbar$ and $\hbar'$ on the two sides. Conjugation preserves the dimension of
  the Lie algebra (\lean{LieIso.finrank_lieAlgebra}). For $\hbar \ne 0$ the observables $X^i$,
  $P_i$, $I$ are linearly independent (\lean{linearIndependent_obsN}): the commutator of a
  vanishing combination with $P_k$ is $i\hbar$ times the coefficient of $X^k$, and with $X^k$
  it is $-i\hbar$ times the coefficient of $P_k$. Their commutators are multiples of $I$, so
  their Lie algebra is their linear span, of dimension $2n + 1$
  (\lean{finrank_lieAlgebra_obsN}). For unitary equivalences the claim fails; this is not
  formalized. The Hermite functions $h_\alpha$
  ($\alpha \in \mathbb{N}^n$) form an orthonormal basis of $L^2(\mathbb{R}^n)$, and $f$ lies in
  $\mathcal{S}(\mathbb{R}^n)$ exactly when its coefficients $c_\alpha$ satisfy
  $\sum_\alpha |c_\alpha|^2 (1 + |\alpha|)^{2k} < \infty$ for all $k$, these sums giving the
  topology (Reed and Simon, \emph{Methods of Modern Mathematical Physics~I}, Thm.~V.13, for
  $n = 1$; B.~Simon, \emph{J. Math. Phys.} 12 (1971) 140). Listing the $\alpha$ by total degree,
  the $j$-th one has $|\alpha| \sim j^{1/n}$, so both weight families define the space $s$ of
  rapidly decreasing sequences. Hence $\mathcal{S}(\mathbb{R}^n)$ and
  $\mathcal{S}(\mathbb{R}^m)$ are unitarily equivalent for all $n, m \ge 1$.
\item Not formalized: [B]STDIST in Theorem~5 and the other remarks of section~5, which are
  informal.
\end{itemize}

\paragraph{Problems found in the brief.} One claim is false, two are imprecise, and two steps
are left implicit.
\begin{itemize}
\item Sec.~5 says that ``separability/second countability would allow us to find a countable
  subset of observables from a bigger set that generates the same topology''. This holds for
  second countability (first countability suffices), but not for separability: see the
  counterexample under ``Differences''.
\item Sec.~5 says that ``if the pure states are second countable, we can always extend the
  topology to the vector space to preserve second countability''. It does not say which
  topology the pure states carry or what extending means. With the quotient topology, no choice
  of topology on $V$ is involved: the topology of Definition~1 is second countable if and only
  if the pure states are.
\item Example~3 says that $\mathcal{S}(\mathbb{R}^n)$ and $\mathcal{S}(\mathbb{R}^m)$ ``are
  isomorphic as quantum vector spaces if and only if $n = m$'', but the brief does not define
  these isomorphisms. The argument of the Remark, that ``no linear transformation can change the
  number of defining observables'', is not valid: the number of defining observables is not an
  invariant. A unitary equivalence $\mathcal{S}(\mathbb{R}^n) \to \mathcal{S}(\mathbb{R})$ carries
  the $2n + 1$ defining observables $X^i$, $P_i$, $I$ to defining observables on
  $\mathcal{S}(\mathbb{R})$ that give the same topology. The claim holds for Lie isomorphisms
  (\lean{nonempty_lieIso_schwartzN_iff}) and fails for unitary equivalences (see
  ``Differences'').
\item The Remark after Example~3 says the $L^2$ seminorms and the sup seminorms dominate each
  other. But the seminorms of Corollary~2 are $\|Af\|$ for every $A$ in the algebra, including
  non-normally-ordered products such as $PX$. Reducing them to $\|x^a \partial^b f\|$ needs the
  commutation relation (normal ordering), which the brief does not mention.
\item Example~3 needs $\hbar \ne 0$: for $\hbar = 0$ we get $P = 0$, and the topology is just
  that of $L^2$.
\end{itemize}

\end{document}
